% ============================================================
%  SCHOLA DOMESTICA — Level 1, Week 4, Lesson 18
%  Topic: Ordering Numbers 0–10 (least to greatest & greatest to least)
%  10 problems | No speed drill
%  Compile with XeLaTeX
% ============================================================
\documentclass[12pt,letterpaper]{article}
\usepackage{sd_style}

\renewcommand{\lessonnum}{18}
\renewcommand{\lessonweek}{4}

\begin{document}

\lessonheader{4}{18}{Ordering Numbers 0–10}

\begin{instructbox}
\noindent\textbf{\color{sd-blue}Instruction} \textit{(read aloud):}\\[4pt]
We can put numbers in order from \textbf{least to greatest} (small to big)
or from \textbf{greatest to least} (big to small).
The number line always goes least to greatest from left to right.
\end{instructbox}

\vspace{6pt}
\begin{center}
\begin{tikzpicture}
  \draw[sd-blue, very thick, ->] (-0.3,0) -- (11.3,0);
  \foreach \n in {0,...,10}{
    \draw[sd-blue, thick] (\n, -0.15) -- (\n, 0.15);
    \node[below] at (\n, -0.25) {\textbf{\n}};
  }
  \node[above] at (2, 0.4) {\color{sd-gold}\small least};
  \node[above] at (9, 0.4) {\color{sd-gold}\small greatest};
  \draw[sd-gold, ->] (2.5,0.35) -- (8.5,0.35);
\end{tikzpicture}
\end{center}

\goldrule
\noindent\textbf{\color{sd-blue}Practice}
\vspace{6pt}

\prob Order from \textbf{least to greatest}.\\[6pt]
\hspace{1cm}{\Large 5 \quad 2 \quad 8} \quad$\longrightarrow$\quad
\ansblank\quad\ansblank\quad\ansblank

\vspace{10pt}
\prob Order from \textbf{least to greatest}.\\[6pt]
\hspace{1cm}{\Large 10 \quad 3 \quad 7 \quad 1} \quad$\longrightarrow$\quad
\ansblank\quad\ansblank\quad\ansblank\quad\ansblank

\vspace{10pt}
\prob Order from \textbf{greatest to least}.\\[6pt]
\hspace{1cm}{\Large 4 \quad 9 \quad 6} \quad$\longrightarrow$\quad
\ansblank\quad\ansblank\quad\ansblank

\vspace{10pt}
\prob Order from \textbf{greatest to least}.\\[6pt]
\hspace{1cm}{\Large 0 \quad 5 \quad 8 \quad 3} \quad$\longrightarrow$\quad
\ansblank\quad\ansblank\quad\ansblank\quad\ansblank

\vspace{10pt}
\prob Write these numbers in order from least to greatest:\\[6pt]
\hspace{1cm}{\Large 9 \quad 0 \quad 6 \quad 4 \quad 2 \quad 7} \quad$\longrightarrow$\quad
\ansblank\quad\ansblank\quad\ansblank\quad\ansblank\quad\ansblank\quad\ansblank

\vspace{10pt}
\prob Which set is in correct order least to greatest? Circle it.\\[6pt]
\hspace{1cm}\textbf{A:} 3, 7, 5, 9 \hspace{2cm} \textbf{B:} 2, 4, 6, 8

\vspace{10pt}
\prob \textit{(Teacher reads)} Five children measured their bean plants.
The heights were 3, 7, 1, 9, and 5 centimeters.
Write them in order from shortest to tallest:\\[6pt]
\hspace{1cm}\ansblank\quad\ansblank\quad\ansblank\quad\ansblank\quad\ansblank

\vspace{10pt}
\prob Fill in the missing numbers to complete the sequence least to greatest:\\[6pt]
\hspace{1cm}2, \ansblank, 6, \ansblank, 10

\vspace{10pt}
\prob Fill in the missing numbers to complete the sequence greatest to least:\\[6pt]
\hspace{1cm}9, \ansblank, 7, \ansblank, 5

\vspace{10pt}
\prob Write any five numbers between 0 and 10, then put them in order.\\[6pt]
My numbers: \ansblank\quad\ansblank\quad\ansblank\quad\ansblank\quad\ansblank\\[6pt]
In order: \ansblank\quad\ansblank\quad\ansblank\quad\ansblank\quad\ansblank

\goldrule

\begin{checkbox}
\noindent\textbf{\color{sd-green}Knowledge Check}\\[4pt]
\setcounter{probnum}{0}
\prob Order least to greatest: 8, 3, 6 \quad$\longrightarrow$\quad\ansblank\quad\ansblank\quad\ansblank\\[8pt]
\prob Order greatest to least: 1, 7, 4 \quad$\longrightarrow$\quad\ansblank\quad\ansblank\quad\ansblank
\end{checkbox}

\vfill
\noindent{\color{sd-blue}$\bigstar$ \textbf{Enrichment — Mayan Numbers 6–10:}}\\[4pt]
The Maya used dots ($\bullet$ = 1) and bars (\rule{0.7cm}{2pt} = 5).\\[4pt]
\begin{tabular}{ccccc}
\rule{0.7cm}{2pt}$\bullet$ &
\rule{0.7cm}{2pt}$\bullet\bullet$ &
\rule{0.7cm}{2pt}$\bullet\bullet\bullet$ &
\rule{0.7cm}{2pt}$\bullet\bullet\bullet\bullet$ &
\rule{0.7cm}{2pt}\rule{0.7cm}{2pt}\\[2pt]
6 & 7 & 8 & 9 & 10
\end{tabular}\\[4pt]
Order these Mayan numbers from least to greatest and write the Hindu-Arabic numerals below.

\vspace{4pt}\hrule\vspace{3pt}
{\footnotesize\color{gray}\textit{Teacher note: Ordering numbers is a prerequisite for understanding
  the number line as a measurement tool. Mastery target: student can correctly
  order any set of numbers 0–10 in either direction without the number line.}}

\end{document}
